\documentclass[11pt]{article}
\usepackage[T1]{fontenc}
\usepackage{lmodern,amsmath,amssymb,amsthm,booktabs,longtable}
\usepackage[margin=1in]{geometry}
\usepackage[colorlinks=true,linkcolor=blue,urlcolor=blue,citecolor=blue]{hyperref}
\usepackage{microtype,xurl}
\newtheorem{theorem}{Theorem}
\newtheorem{lemma}[theorem]{Lemma}
\newtheorem{proposition}[theorem]{Proposition}
\newcommand{\Z}{\mathbb Z}
\newcommand{\Q}{\mathbb Q}
\newcommand{\OK}{\mathcal O_K}
\newcommand{\pos}[1]{\left(#1\right)_+}
\newcommand{\Tn}{\mathcal T}
\title{A smoothed Frobenius determinant bound\\for the tenth discriminant\\
\large Version 0.6.0: \(q<570\)}
\author{Research note prepared with Codex}
\date{20 October 2026 --- independent mathematical review pending}
\begin{document}
\maketitle

\begin{abstract}
We obtain \(q<570\) for an imaginary quadratic field of class number
one. The new argument excludes \(570\le q\le2000\); the previously
audited v5 result supplies the range above 2000. The proof retains a
smoothed prime sum already present in the determinant valuation,
uses it directly in an acyclic factorial refinement, and adds compatible
higher Frobenius factors. A pointwise inequality couples derivative
orders with interpolation radii. A positive kernel formed from a squared
polynomial completes the comparison with the quadratic norm series.
All determinant parameters are fixed, apart from the scale and its
auxiliary-prime interval. The numerical certificate uses rational
intervals and Python's standard library. The external number-theoretic
inputs are quadratic ideal theory and explicit ordinary-prime estimates
of Rosser and Schoenfeld. No quasi-Riemann hypothesis or zero-free theorem
for quadratic Dirichlet functions is used.
\end{abstract}

\section{Scope, provenance, and the result}
For a negative fundamental discriminant \(D\), put \(q=|D|\) and
\(K=\Q(\sqrt D)\). We address Tao's request for tenth-discriminant
bounds obtained without the stronger zero-location machinery discussed
in the comments to his challenge \cite{tao,taocomment}.

\begin{theorem}\label{main}
If \(h(K)=1\), then \(q<570\).
\end{theorem}

The sole earlier-version input is the independently audited implication
\(h(K)=1\Rightarrow q<2000\) from v5 \cite{v5}. The entire new exclusion
of \([570,2000]\) is proved here, relative to the classical prime estimates
stated below. None of the intermediate unpublished 1600, 1100, or 1050
drafts is a dependency. Discriminants in the new interval are not
enumerated to obtain the analytic cutoff.

The framework is Green's derivative determinant \cite{green}, with
van Dijk's auxiliary inert prime for nonvanishing \cite{vandijk}.
Green's comment also records a related unpublished argument of Jacob
Pescador \cite{greencomment}; we have not used that unpublished text.
Pratt proposed combining auxiliary inert primes from quadratic norm
counts, Green's stronger divisibility, and disk interpolation
\cite{pratt}. Kerr's suggestion to retain higher Frobenius relations,
including the ninth-power relation at 3, also informed the valuation
refinements \cite{kerrcomment}. The smoothed comparison, direct factorial
refinement, compatible residual factors, and joint archimedean estimate
below were developed in this AI-assisted research session. No priority
claim is made for these refinements.
The Poisson expansion of the quadratic norm series is classical;
Ross and Tao discussed its usefulness in this challenge
\cite{ross,taocomment}. We use it in the half-plane \(s>1\), with
an explicit error bound for the signed kernel comparison.

This v6 draft has not yet had an independent mathematical audit. Exact
computational checks are distinct from an audit of the general argument
and from formal verification. Earlier versions retain their own review
records; those reviews do not cover this draft.

\section{Elementary consequences of class number one}
Assume henceforth that \(h(K)=1\) and \(570\le q\le2000\).
Every nonrational algebraic integer has norm at least \(q/4\).
A noninert rational prime has an integral ideal of norm \(p\);
under class number one its generator is nonrational. Consequently
\begin{equation}\label{small-inert}
 p<q/4\quad\Longrightarrow\quad p\text{ is inert}.
\end{equation}
In particular \(q\) is an odd prime: an even discriminant has a small
ramified prime 2, and a composite odd discriminant has a ramified prime
at most \(\sqrt q<q/4\). Inertness of 2 gives \(q\equiv3\pmod8\).
The only units are \(\pm1\), and
\[
 \OK=\Z+\Z\omega,\qquad \omega=(1+i\sqrt q)/2,\qquad
 Q_q(a,b)=|a+b\omega|^2=(a+b/2)^2+qb^2/4.
\]
Write
\[
 v=\sqrt q/2,\qquad \rho=\tfrac12\sqrt{1+v^2}.
\]
The lattice has covolume \(v\) and a rectangular tiling of covering
radius \(\rho\). Its disk count \(N(r)\) therefore satisfies
\begin{equation}\label{disk}
 \frac{\pi\max(r-\rho,0)^2}{v}\le N(r)
 \le\frac{\pi(r+\rho)^2}{v}.
\end{equation}
Every principal ideal of norm \(m\) is a rotation and dilation by
\(\sqrt m\) of this lattice.

\section{The ordinary-prime inputs}
Put \(\vartheta(x)=\sum_{p\le x}\log p\) and
\(A(x)=\sum_{p\le x}\log p/p\). We use Rosser--Schoenfeld
\cite{rs}, with \(c_0=1.01624\):
\begin{align}
 \vartheta(x)&<c_0x,\label{theta-global}\\
 x\left(1-\frac1{2\log x}\right)&<\vartheta(x)
       \quad(x\ge563),\label{theta-lower}\\
 \vartheta(x)&<x\left(1+\frac1{2\log x}\right)\quad(x>1),\label{theta-upper}\\
 A(x)&>\log x-1.334-\frac1{2\log x}\quad(x\ge319).
 \label{mertens}
\end{align}
The last constant follows from their Theorem 6 and the explicit prime
sum at 319: the certificate checks
\(\log319-A(319)+1/(2\log319)<1.334\).

For sharper bounds at our scales, Table I on printed page 90 gives
pairs \((b,\epsilon_b)\) such that
\((1-\epsilon_b)x<\psi(x)<(1+\epsilon_b)x\) for \(x\ge e^b\).
The corollary to their Theorem 14 gives
\(0\le\psi(x)-\vartheta(x)<1.02\sqrt x+3x^{1/3}\).
Thus, uniformly for \(\log x\ge m\ge b\),
\begin{equation}\label{tabletheta}
 1-\epsilon_b-1.02e^{-m/2}-3e^{-2m/3}
 <\frac{\vartheta(x)}x<1+\epsilon_b.
\end{equation}
Let \(c(m),d(m)\) be any certified upper and lower bounds obtained
by taking the best of these inequalities and
\eqref{theta-global}--\eqref{theta-upper}. The table entries used are
included as exact decimals in \texttt{parameters.json}.

We also use \(A(R)<2\log R\) at our scales. This follows without a
further prime estimate: \(\log(R!)\ge RA(R)-\vartheta(R)\), so
\(A(R)\le\log R+c_0<2\log R\).

\section{A sieve providing the auxiliary inert prime}\label{sieve}
Fix one band \(Q_0\le q\le Q_1\). Form the fixed, finite list of pairs
\begin{equation}\label{normlist}
 b>0,\quad 2a+b>0,\quad \gcd(a,b)=1,\quad
 Q_{Q_1}(a,b)<(Q_0/4)^2,
\end{equation}
ordered by \((Q_{Q_1}(a,b),a,b)\). Write \(p_j(q)=Q_q(a_j,b_j)\),
\(p_j^-=Q_{Q_0}(a_j,b_j)\), and \(p_j^+=Q_{Q_1}(a_j,b_j)\).

Each \(p_j(q)\) is prime. A composite norm below \((q/4)^2\) has a
prime factor below \(q/4\); inertness would force that prime to divide
both coordinates. It is split rather than ramified: generators of the
unique prime ideal above \(q\) have zero real part. These primes are
distinct, because the four generators of a split prime norm are
\(\alpha,-\alpha,\bar\alpha,-\bar\alpha\), and the chosen quadrant
contains exactly one of them.

Choose the first \(n\) such split primes, and the first \(h\) rational
primes \(r>7\) with \(r<Q_0/4\). Sieve by the two prime ideals above
each selected \(p_j\), and by each ideal \(r\OK\). Their norms are
\(p_j,p_j,r^2\), and they are pairwise coprime. Also impose the exact
condition \(\gcd(a,b,210)=1\).

If an element belongs to \(k\) of the selected ideals, then
\(1-k+\binom{k}{2}\) is an upper bound for the indicator that it
belongs to none. Apply this second Bonferroni inequality, multiplied
by the exact 210 indicator, to the annulus
\(U<|\alpha|^2\le\kappa U\). A split prime in this interval has four
generators. Assume \(U>Q_1\) and \(U>\max_{j\le n}p_j^+\), so the
sieving primes themselves are not counted.

For an ideal of norm \(m\), the disk estimate gives an annulus count
with area term \(\pi(\kappa-1)U/(vm)\), boundary error at most
\(2\pi\rho(1+\sqrt\kappa)\sqrt U/(v\sqrt m)\), and constant error
at most \(2\pi\rho^2/v\). The 210 expansion has
\[
 \sum_{d\mid210}\frac{\mu(d)}{d^2}=\frac{768}{1225},\qquad
 \sum_{d\mid210}\frac{|\mu(d)|}{d}=\frac{96}{35},\qquad
 \sum_{d\mid210}|\mu(d)|=16.
\]
Define
\begin{align*}
 u&=2\sum_{j\le n}(p_j^+)^{-1}+\sum_{r}r^{-2},&
 u_2&=2\sum_{j\le n}(p_j^+)^{-2}+\sum_r r^{-4},\\
 z&=2\sum_{j\le n}(p_j^-)^{-1/2}+\sum_r r^{-1},&
 z_2&=2\sum_{j\le n}(p_j^-)^{-1}+\sum_r r^{-2},
 \qquad m_0=2n+h.
\end{align*}
The certificate checks that the analogous \(u\) at the minus endpoints
is less than 1. The area polynomial then decreases in every reciprocal
ideal norm, since its partial derivative is \(-1\) plus the sum of
the other reciprocals. Therefore uniform upper bounds for the number
of split primes in \((U,\kappa U]\) are
\begin{equation}\label{split-annulus}
 A_*(\kappa-1)U+B_*(1+\sqrt\kappa)\sqrt U+C_*,
\end{equation}
where \(v_-=\sqrt{Q_0}/2\), \(v_+=\sqrt{Q_1}/2\), and
\begin{align*}
 A_*&=\frac{\pi}{4v_-}\frac{768}{1225}
          \left(1-u+\frac{u^2-u_2}{2}\right),\\
 B_*&=\frac\pi4\sqrt{1+v_-^{-2}}\frac{96}{35}
          \left(1+z+\frac{z^2-z_2}{2}\right),\\
 C_*&=2\pi(v_++v_+^{-1})
          \left(1+\frac{m_0(m_0+1)}2\right).
\end{align*}

For a parameter row, put \(M=2^a\), and let
\(2^a\le U\le2^b\). Set \(d=d(\log M)\), \(c=c(\log M)\).
An inert prime exists in \((U,\kappa U]\) if
\begin{equation}\label{sieve-margin}
 G(x)=\kappa d-c-(x+\log\kappa)
 \{A_*(\kappa-1)+B_*(1+\sqrt\kappa)e^{-x/2}+C_*e^{-x}\}>0,
 \quad x=\log U.
\end{equation}
Indeed, the total prime logarithmic weight exceeds the possible split
weight, and the ramified prime is below the interval. Since
\(x+\log\kappa>4\), the function \(G\) is concave. It suffices to
check its two endpoints. These are exact rational interval checks in
the certificate; all \(\kappa,n,h,a,b\) are in the parameter file.

\section{Nonvanishing with the prime as monomial width}\label{rank}
Take an inert prime \(p_0\in(U,\kappa U]\), and set \(R=p_0\).
Throughout the new proof use the fixed constants
\begin{equation}\label{fixed}
 t=\frac3{64},\qquad C=\frac{27}{64},\qquad
 S=\lfloor tR\rfloor,\qquad J=\lceil R/9\rceil,
 \qquad T=\sqrt{CRv/\pi}+\rho.
\end{equation}
Use all \(n\) lattice points \(z_i\) in the disk of radius \(T\).
Then \(n\ge CR\), so \(nJ\ge tR^2\ge RS\).
The matrix has rows \(0\le r<R,0\le s<S\), columns \((i,k)\),
\(0\le k<J\), and entries
\[
 \partial_z^k(z^r w^s)\big|_{(z_i,\bar z_i)}.
\]

\begin{lemma}\label{nonzero}
If \(2T<R\) and \(J\le R\), this matrix has full row rank.
\end{lemma}
\begin{proof}
Reduce modulo \(p_0\OK\), a field of \(p_0^2\) elements. The points
are distinct: a nonzero element of \(p_0\OK\) has modulus at least
\(p_0\), whereas two points in the disk differ by less than \(2T\).
Frobenius gives \(\bar z_i=z_i^{p_0}\).
If a nonzero row combination vanished, substitute \(w=z^{p_0}\) in
its polynomial. The exponents \(r+p_0s\) are distinct, so the result
is a nonzero polynomial of degree at most \(RS-1\). For \(k<J\le p_0\),
its ordinary derivatives agree with the specified \(z\)-derivatives
in characteristic \(p_0\). Each node is thus a zero of multiplicity
at least \(J\), since \(k!\ne0\) for \(k<p_0\). This gives at least
\(nJ\ge RS\) zeros, a contradiction. Full rank after reduction
implies full rank in characteristic zero.
\end{proof}

Fix a nonzero \(RS\)-square minor \(\mathcal D\), and put
\(N=RS\) and \(\Delta=|\mathcal D|^2\in\Z_{>0}\).
Different triangular row changes below may be used for different
primes: they leave this same determinant unchanged.

\section{The smoothed prime sum already in the valuation}
For \(0<u\le1\), put
\begin{equation}\label{fdef}
 f(u)=2au-a(a+1)u^2,\qquad a=\lfloor1/u\rfloor;
 \qquad f(u)=0\ (u>1),\quad f(0)=1.
\end{equation}
This is continuous, nonnegative, nonincreasing, and at most 1. Also
\begin{equation}\label{fidentity}
 f(u)=2u\sum_{m\ge1}(1-mu)_+,\qquad
 \frac2{R^2}\sum_{r=0}^{R-1}\left\lfloor\frac rp\right\rfloor
       =\frac{f(p/R)}p.
\end{equation}
Define the noninert sum
\begin{equation}\label{Tdef}
 \Tn(X)=\sum_{p\text{ noninert}}\frac{\log p}{p}f(p/X).
\end{equation}
It is nonnegative and nondecreasing in \(X\), and is at most the
ordinary cumulative noninert sum.

Green's triangular change, using
\(a_{r,s}=\min(\lfloor r/p\rfloor,S-1-s)\), replaces \(z^rw^s\) by
\(z^{r-pa_{r,s}}(z^p-w)^{a_{r,s}}w^s\). For inert \(p\), every
ordinary derivative of \(z^p-w\) is divisible by \(p\), so
\[
 v_p(\Delta)\ge2\sum_{r,s}a_{r,s}.
\]
Write this normalized valuation as \(e(R,S,p)\). Counting layers gives
\begin{equation}\label{elayers}
 e(R,S,p)=\frac2{R^2S}\sum_{h=1}^{S-1}(S-h)(R-ph)_+.
\end{equation}
For \(p(S-1)\le R-1\), this equals
\((S-1)/R-p(S^2-1)/(3R^2)\). In general,
\begin{equation}\label{etail}
 e(R,S,p)\ge\frac{f(p/R)}p
       -\frac{R-1}{2RSp}
       -\frac{(R-1)(2R-1)}{6RSp^2}.
\end{equation}
Indeed, sum
\(\sum_{s<S}\min(k,S-1-s)\ge Sk-k(k+1)/2\), then use
\(k(k+1)\le r^2/p^2+r/p\) and \eqref{fidentity}.

The main term in the inert-prime sum in \eqref{etail} is exactly
\[
 \sum_p\frac{\log p}{p}f(p/R)-\Tn(R).
\]
Keeping \(\Tn\), instead of replacing it by the larger ordinary
cumulative sum, is an important improvement.

\subsection{The all-prime floor remainder}
If \(R=ap+b\), the exact floor identity gives
\[
 \frac{f(p/R)}p=\frac1p-\frac1R+\frac1R\delta_a(p/R),\qquad
 \delta_a(u)=2a+1-u^{-1}-a(a+1)u.
\]
On \([1/(a+1),1/a]\), \(\delta_a\) is nonnegative, vanishes at
the endpoints, and has one maximum. Partial summation with
\(d_*x\le\vartheta(x)\le c_*x\), valid for \(x\ge M/33\), gives
\begin{equation}\label{Gfloor}
 G_{\rm floor}=\max(0,Bd_*-Ac_*),\quad
 A=\frac{\log33-H_{33}+1}{2},\quad
 B=\frac{H_{32}-\log33}{2}.
\end{equation}
Here \(H_n=\sum_{j\le n}1/j\). The increasing and decreasing pieces
contribute integrals
\(\frac12\log(1+1/a)-\frac1{2(a+1)}\) and
\(\frac1{2a}-\frac12\log(1+1/a)\), respectively; sum for
\(1\le a\le32\). Thus
\begin{equation}\label{allf}
 \sum_p\frac{\log p}{p}f(p/R)
 \ge A(R)-c(\log M)+G_{\rm floor}\qquad(R\ge M).
\end{equation}

\subsection{Removing known noninert primes from the quadratic loss}
Put \(H=\lfloor1/t\rfloor=21\),
\[
 a_H=\sum_{p\le H}\frac{\log p}p,\quad
 \theta_H=\sum_{p\le H}\log p,\quad
 W_H=\sum_{p\le H}\log p\,(t-pt^2/3).
\]
Let \(B_H\) be a certified upper bound for
\(\sum_{p>H}\log p/p^2\). Specifically, with \(P=10000\),
partial summation and \eqref{theta-global} give
\[
 B_H=\sum_{H<p\le P}\frac{\log p}{p^2}
               +\frac{2c_0}{P}-\frac{\vartheta(P)}{P^2}
 >\sum_{p>H}\frac{\log p}{p^2}.
\]
With all pairs in \eqref{normlist}, define
\begin{equation}\label{Ktwo}
 K_2=\sum_j\frac{\log p_j^+}{(p_j^+)^2}
             +\frac{\log Q_1}{Q_1^2}.
\end{equation}
The summands decrease with the prime, so the inert quadratic tail
is at most \(B_H-K_2\). Combining \eqref{etail}, \eqref{allf}, the
exact small-prime formula, and \(A(R)<2\log R\), gives
\begin{align}\label{basevaluation}
 \frac{\log\Delta}{NR}\ge{}&
 A(R)-c(\log M)+G_{\rm floor}-\Tn(R)-a_H+W_H\\
 &-\frac{B_H-K_2}{3(t-1/M)}
 -\frac{2\theta_H}{M}-\frac{\log M}{(t-1/M)M}.
 \notag
\end{align}
The next two sections give compatible additional gains.

\section{Compatible higher Frobenius factors}\label{frobenius}
At an inert prime use
\[
 G=z^p-w,\qquad H_1=z^{p^2}-z,\qquad H_2=w^{p^2}-w.
\]
For each row let
\[
 a=\min(\lfloor r/p\rfloor,S-1-s),\quad
 c=\lfloor(r-pa)/p^2\rfloor,\quad b=\lfloor s/p^2\rfloor,
\]
and replace the row polynomial by
\begin{equation}\label{threefactor}
 z^{r-pa-p^2c}G^aH_1^c\,w^{s-p^2b}H_2^b.
\end{equation}
Its leading term is \(z^rw^s\); all other terms have smaller
lexicographic index and remain in the rectangle, since their
\(z\)-degree is at most \(r\) and their \(w\)-degree at most
\(s+a\le S-1\).

All three factors have values divisible by \(p\). A Hasse derivative
can lose a factor \(G\) only at order \(p\), and a factor \(H_1\)
only at order 1 or \(p^2\). At total order \(k\), the number of lost
factors is therefore at most
\[
 L=\min(k,c)+\left\lfloor\frac{k-\min(k,c)}p\right\rfloor.
\]
Ordinary differentiation supplies \(k!\). The elementary inequality
\[
 c+\lfloor k/p\rfloor-L\ge\lfloor c/p\rfloor
\]
holds both for \(k\ge c\) and \(k<c\). Consequently the entry
valuation is at least
\begin{equation}\label{newval}
 a+b+\lfloor c/p\rfloor+\sum_{e\ge2}\lfloor k/p^e\rfloor.
\end{equation}
The additional residual row factor is
\(\lfloor(r-pa)/p^3\rfloor\). Moreover, beyond this bound one may
retain
\begin{equation}\label{coupled-entry}
 \max\left(0,(1-1/p)(c-k)-1\right).
\end{equation}
For \(k<c\), this follows by bounding \(\lfloor k/p\rfloor\)
below by \(k/p-1\); for \(k\ge c\), \eqref{newval} suffices.

\subsection{Explicit gains used by the certificate}
Suppose \(R\ge M\), put \(\varepsilon=1/M\), and choose
\(b_0\le S/n\). At most \(n\) columns have each derivative order.
Their mean order is at least \((b_0R-1)/2\), and at most \(j_*R\),
where
\begin{equation}\label{jstar}
 j=1/9+\varepsilon,\qquad j_*=j-b_0/2+b_0/M.
\end{equation}
The latter is a conservative consequence of sorting the orders:
the \(i\)th is at most \(J-1-\lfloor(N-1-i)/n\rfloor\).

The \(b\) and higher-column terms in \eqref{newval} give
\begin{equation}\label{Gsmall}
 G_{\rm small}=\sum_{p<Q_0/4}\log p\left\{
 \sum_{\substack{e\ge2\\p^e<M}}
       \left(\frac{b_0-\varepsilon}{p^e}-2\varepsilon\right)_+
 +\left(\frac{t-\varepsilon}{p^2}-3\varepsilon\right)_+
 \right\}.
\end{equation}
For the residual factor, define
\[
 A_p(t)=
 \begin{cases}
  1/(2p^2)-t/(2p)+t^2/6,&pt\le1,\\
  1/(6tp^3),&pt\ge1.
 \end{cases}
\]
Then the further gains are
\begin{align}
 G_{\rm res}&=\sum_{p<Q_0/4}\log p
       \left(\frac{2A_p(t)}p-(2+2/p^3)\varepsilon\right)_+,
       \label{Gres}\\
 G_{\rm coup}&=2\sum_{p<Q_0/4}\log p
       \left((1-1/p)\{A_p(t)-(1+1/p^2)\varepsilon-j_*\}
                         -\varepsilon\right)_+.
       \label{Gcoup}
\end{align}
Here is a direct integral justification. Put \(x_r=r/R\) and
\(y_s=(S-1-s)/R\). We have
\(c/R\ge (x_r-py_s)_+/p^2-1/R\).
On the cell \([x_r,x_r+1/R]\times[y_s,y_s+1/R]\),
\((x-py)_+\le(x_r-py_s)_++1/R\).
The average of \((x-py)_+/p^2\) over
\([0,1]\times[0,S/R]\) equals \(A_p(S/R)\), which decreases
as \(S/R\) increases. Thus the normalized mean of
\(c\) is at least \(A_p(t)-(1+1/p^2)/M\).
For \eqref{Gres}, use the exact identity
\(\lfloor c/p\rfloor=\lfloor(r-pa)/p^3\rfloor\).
The same cell argument, now divided by \(p^3\), bounds its
normalized mean below by \(A_p(t)/p-(1+1/p^3)/M\).
The mean of a positive part is at least the positive
part of its mean, so \eqref{coupled-entry} and \eqref{jstar} give
\eqref{Gcoup}. All these terms are added to \eqref{basevaluation}.

\section{The factorial refinement uses the same smoothed sum}
At a split prime, replace \(w^s\) by its falling factorial
\((w)_s\). Since \(\OK/p^e\OK\simeq(\Z/p^e\Z)^2\),
\(p^{v_p(s!)}\) divides its value. The column factor \(k!\) remains.
At the ramified prime keep the column factorial only. These row
changes are triangular, and therefore bound the same determinant.

Suppose a nonnegative nondecreasing function \(\Phi\) is already
known to satisfy \(\Tn(e^x)\ge\Phi(x)\) throughout the band.
The first-power part of the split row factorial contributes exactly
\((S/R)\Tn_{\rm split}(S)\), by \eqref{fidentity}. It is therefore
at least
\[
 (t-\varepsilon)
 \left(\Phi(\log(M(t-\varepsilon)))-\frac{\log Q_0}{Q_0}\right)_+.
\]
For the columns put \(K_0=\lfloor b_0R\rfloor\). The increasing
chosen orders are at least \(\lfloor i/n\rfloor\).
Since \(N/n\ge K_0\), the average of \(\lfloor k/p\rfloor\)
over these columns is at least its average for \(0\le k<K_0\).
Using \eqref{fidentity} again gives a contribution at least
\((K_0/R)\Tn(K_0)\). Hence a compatible total gain is
\begin{align}\label{Gfeedback}
 G_{\rm fb}={}&(b_0-\varepsilon)
       \Phi(\log(M(b_0-\varepsilon)))\\
 &+(t-\varepsilon)
       \left(\Phi(\log(M(t-\varepsilon)))-\frac{\log Q_0}{Q_0}\right)_+.
 \notag
\end{align}
There is no fractional-part approximation in this step. We start
with \(\Phi=0\), compute a new lower bound, and repeat eight times.
Every stage uses only the preceding stage, so the procedure is acyclic.

\section{A joint bound for derivative orders and radii}\label{arch}
We give the pointwise inequality underlying the archimedean estimate.
Fix \(\tau=3/10\), and, for \(0\le k\le u\le1\), put
\[
 F(u,k)=u\log u-(u-k)\log(u-k)-k,\qquad
 \phi(u)=\tau(u\log u-u).
\]
At zero use the continuous values. For \(0<a\le1\), define
\(\psi(k,a)=\max_{k\le u\le1}\{(u-k)\log a+F(u,k)-\phi(u)\}\).
The maximized function is strictly concave. If \(a\le1-k\), its
critical point satisfies \(u-k=au^{1-\tau}\), and
\[
 \psi(k,a)=\tau u+(1-\tau)k\log u-k.
\]
If \(a\ge1-k\), the maximum is at 1 and
\[
 \psi(k,a)=(1-k)\log\frac a{1-k}-k+\tau.
\]
Differentiation shows that \(\psi\) decreases in \(k\), increases in
\(a\), and is 1-Lipschitz in \(a\). In particular, for \(k\le j\),
\[
 \psi(k,a)+B\ge0,\qquad B=(1-\tau)j(1-\log j).
\]

Let \(I(k)=\int_0^1 2a\psi(k,a)\,da\) and
\(K(j)=\int_0^j I(k)\,dk\). We have \(I(0)=2\tau^2/(2\tau+1)\).
Changing variables from \(a\) to the critical point gives
\[
 I'(k)=-(1-\tau)\int_k^1(u-k)^2u^{2\tau-3}\,du
               +\log(1-k)+k-k^2/2.
\]
Thus, writing \(A=2\tau\),
\begin{align}\label{Kformula}
 K(j)={}&j\frac{2\tau^2}{2\tau+1}
 -(1-\tau)\left\{\frac{j^2}{2A}-\frac{j^3}{3(A-1)}
 +\frac{j^4}{12(A-2)}
 -\frac{2j^{A+2}}{A(A-1)(A-2)(A+1)(A+2)}\right\}\\
 &+\frac j2-\frac{3j^2}4+\frac{(1-j)^2}2\log(1-j)
        +\frac{j^3}6-\frac{j^4}{24}.
 \notag
\end{align}

For an actual row and derivative order, take \(u=(r+1)/R\) and
\(k'=k/R\). Increasing-log integration gives
\(\log(r)_k\le k\log R+RF(u,k')\).
Clip the normalized radius to \(a=\max(1,|z|)/T\). A nonzero entry
satisfies
\begin{align*}
 \log|\partial_z^k(z^rw^s)|\le{}&
 (r+s)\log T+k\log(R/T)\\
 &+R\{\phi(u)+\psi(k',a)\}+s\log a-\log a.
\end{align*}
The final term is at most \(\log T\). Since \(\phi\) decreases, its
right Riemann sum gives a normalized row contribution at most
\(-3\tau/2\) to \(\log\Delta/(NR)\).

For uniform estimates set \(\varepsilon=1/M\),
\begin{equation}\label{geometry}
 r_* =\sqrt{CMv_+/\pi},\quad e\ge\rho_+/r_*,\quad
 \nu=C(1+2e)^2,\quad b_0\le(t-\varepsilon)/\nu,
 \quad g=\tfrac12\log(C/(2\pi))+\log(1+e).
\end{equation}
Here \(v_+=\sqrt{Q_1}/2\), \(\rho_+=\sqrt{1+v_+^2}/2\).
Then \(n\le\nu R\), \(b_0\le S/n\), and
\(\log T\le\frac12\log R+\frac14\log q+g\).
The monotonicity of \(\rho/\sqrt v=\frac12\sqrt{v+v^{-1}}\)
justifies use of the upper band endpoint.

Lattice-cell integration, using nonnegativity and the Lipschitz property,
gives for each order
\[
 \sum_{|z|\le T}\{\psi(k/R,a_z)+B\}
 \le\nu R\left\{I(k/R)+B+\frac{5e}{3(1+e)}\right\}.
\]
Indeed, cells lie in the disk of radius \(T+\rho\). Its normalized
radial argument is \(((1+2e)a+e)/(1+e)\), whose difference from
\(a\) is \(e(1+a)/(1+e)\). Integrating against \(2a\,da\) gives
the displayed error. Clipping at 1 only improves the estimate;
\(\rho\ge1\) also covers the clipped origin.

We may fill in unchosen columns before subtracting \(NB\).
Since \(I\) decreases and \(I+B\ge0\), the order sum is at most
\[
 R\left\{K(j)+j\left(B+\frac{5e}{3(1+e)}\right)\right\}+I(0)+B.
\]
It follows that the normalized joint contribution is at most
\begin{equation}\label{Djoint}
 D_*=-\frac{3\tau}2+
 \frac{2\nu}{t-\varepsilon}
 \left\{K(j)+j\left(B+\frac{5e}{3(1+e)}\right)
                  +\varepsilon(I(0)+B)\right\}-2B.
\end{equation}

\subsection{The unused \texorpdfstring{\(s\)}{s}-radius contribution}
At least \(\lambda(a)N\) chosen columns have radius at most \(aT\), where
\[
 \lambda(a)=\left(1-\frac{jC}{t-\varepsilon}
   \{(1+2e)^2-\max(a(1+e)-e,0)^2\}\right)_+.
\]
Among that many distinct rows the normalized sum of the \(s\)'s is
at least \(((t-\varepsilon)\lambda(a)^2-\varepsilon)/2\).
Put
\[
 \delta=\frac{jC}{t-\varepsilon}
       \left\{(1+2e)^2+\frac e{2+e}\right\}-1.
\]
Since \(\max(a(1+e)-e,0)^2\ge a^2-e/(2+e)\) and
\(jC/(t-\varepsilon)\ge1\), we have
\(\lambda(a)\ge(a^2-\delta)_+\).
The certificate checks \(0<\delta<1\) and \(\delta\ge e\), so
\(\sqrt\delta\ge1/T\). Integrating from \(\sqrt\delta\) to 1 gives
the additional saving
\begin{equation}\label{Gs}
 G_s=\left((t-\varepsilon)
       \left\{\frac14-\delta+\frac{3\delta^2}4
                    -\frac{\delta^2}2\log\delta\right\}
                     +\frac\varepsilon2\log\delta\right)_+.
\end{equation}

There are at most \(N!\) terms in the determinant. Using
\(\log N!\le N\log N\), \eqref{jstar}, and \(1<T<R\), we obtain
\begin{equation}\label{arch-final}
 \frac{\log\Delta}{NR}\le
 (1-\alpha)\log R+\frac\beta4\log q+\beta g
       +D_*+\frac{2\log M}{M}-G_s+\frac{2\log(tM^2)}M,
\end{equation}
where
\(\alpha=(1-t)/2-j_*\), \(\beta=1+t-2j_*\).
The functions \((\log R)/R\) and \(\log(tR^2)/R\) decrease on the
certified ranges. Notice that the joint estimate replaces the former
separate derivative and radius estimates; incompatible savings are
not added.

\section{The certified affine bounds and finite iteration}\label{affine}
Combine \eqref{basevaluation}, \eqref{mertens},
\eqref{Gsmall}--\eqref{Gcoup}, \eqref{Gfeedback}, and
\eqref{arch-final}. This gives
\begin{equation}\label{Tline}
 \Tn(R)\ge\alpha\log R-\frac\beta4\log q-\mathcal K,
\end{equation}
where the fully explicit constant is
\begin{align*}
 \mathcal K={}&a_H-W_H+c(\log M)
 +\frac{B_H-K_2}{3(t-\varepsilon)}+\beta g+D_*
 +\frac{2\log M}{M}-G_s\\
 &+\frac{2\theta_H}{M}+\frac{\log M}{(t-\varepsilon)M}
 +\frac{2\log(tM^2)}M+1.334+\frac1{2\log M}\\
 &-G_{\rm floor}-G_{\rm small}-G_{\rm res}-G_{\rm coup}-G_{\rm fb}.
\end{align*}
For every \(U\in[2^a,2^b]\), the auxiliary prime satisfies
\(U<R\le\kappa U\). Monotonicity of \(\Tn\) and \(\alpha>0\)
therefore transfer \eqref{Tline} to \(X=\kappa U\), with the
additional intercept cost \(\alpha\log\kappa\).

The certificate rounds slopes down and intercept costs up, obtaining
\[
 \Tn(e^x)\ge A_i x-D_i-\gamma\log(q/Q_0),\qquad
 \gamma=3/10,
\]
on the indicated finite interval. It checks \(\beta/4<\gamma\).
At a finite endpoint the endpoint value remains a lower bound at
all later \(x\). The maximum with zero is nondecreasing.
For each of eight refinement stages, use this preceding maximum
minus \(\gamma\log(Q_1/Q_0)\), clipped at zero, as \(\Phi\) in
\eqref{Gfeedback}. The initial stage uses \(\Phi=0\).

Write \(L(x)\) for a certified piecewise affine selection from the
final bounds at \(Q_0\), including zero and frozen endpoint values.
Then throughout the band,
\begin{equation}\label{Luniform}
 \Tn(e^x)\ge L(x)-\gamma\log(q/Q_0).
\end{equation}
The exact parameter and output files record all lines, stages, and
selected segments. No optimizer is part of the certificate.

\section{Positive kernels and the quadratic norm series}\label{mellin}
Choose positive rational \(d,e\) and a nonzero integer polynomial
\(P\), and normalize
\[
 h(x)=Z^{-1}e^{-dx}P(e^{-ex})^2\quad(x\ge0),\qquad
 Z=\int_0^\infty e^{-dx}P(e^{-ex})^2\,dx.
\]
Then \(h\ge0\) and \(\int h=1\). Expand
\(h(x)=\sum_i c_i d_i e^{-d_i x}\), with \(d_i=d+ie\).
All \(c_i\) and \(Z\) are exact rationals. Integrating
\eqref{Luniform} gives
\begin{equation}\label{weightedlower}
 \int_0^\infty h(x)\Tn(e^x)\,dx
 \ge\sum_i c_i\!\int_0^\infty d_i e^{-d_i x}L(x)\,dx
             -\gamma\log(q/Q_0).
\end{equation}
One and the same \(L\) is used in every signed term.

Define
\[
 W(y)=\int_y^\infty h(x)f(e^{y-x})\,dx.
\]
This is nonnegative, nonincreasing, and at most 1. Tonelli's theorem
and \eqref{fidentity} show that
\begin{equation}\label{multiplier}
 W(y)=\sum_i c_i F(d_i)e^{-d_i y},\qquad
 F(d)=\int_0^1 d\,u^{d-1}f(u)\,du
       =\frac{2d\zeta(1+d)}{(1+d)(2+d)}.
\end{equation}
The right side follows by integrating each
\(2d\,u^d(1-mu)_+\) from 0 to \(1/m\).
Thus the left side of \eqref{weightedlower} is
\(\sum_{p\text{ noninert}}(\log p)/p\,W(\log p)\).

For \(s>1\), put
\[
 \mathcal A_q(s)=\sum_{b\ge1}\sum_{a\in\Z}Q_q(a,b)^{-s},\qquad
 B_q(s)=1+\mathcal A_q(s)/\zeta(2s).
\]
Class number one and the unit group give
\(\zeta_K(s)=\zeta(2s)+\mathcal A_q(s)\). Its Euler product implies
\[
 -\frac{B_q'(s)}{2B_q(s)}
 =\sum_{p\text{ split}}\sum_{\substack{m\ge1\\m\text{ odd}}}
      \frac{\log p}{p^{ms}}
 +\frac{\log q}{2}\sum_{m\ge1}(-1)^{m-1}q^{-ms}.
\]
Multiply by the signed coefficients \(c_iF(d_i)\) and sum at
\(s_i=1+d_i\). Since \(W\ge0\), the discarded higher odd powers
have nonnegative weights. Adding
\(\frac12\log q\sum_{m\ge1}q^{-m}W(m\log q)\) also covers the
ramified first-power term. It is at most
\begin{equation}\label{ramerror}
 \frac{\log Q_0}{2(Q_0-1)}.
\end{equation}

\section{A small explicit norm-series error}\label{poisson}
Set
\[
 C(s)=\sqrt\pi\frac{\Gamma(s-1/2)}{\Gamma(s)},\quad
 g(s)=\psi(s)-\psi(s-1/2),\quad
 \mathcal F_q(s)=C(s)v^{1-2s}\zeta(2s-1).
\]
Poisson summation in the real coordinate gives a particularly small
error for \(\mathcal A_q-\mathcal F_q\). We supply a bound suitable
for the signed kernel coefficients.

Choose any rational \(0<a_0<v_-\). For
\(f_{s,t}(x)=(x^2+t^2)^{-s}\), shift the Fourier contour by
\(y=t-a_0\), where \(t=bv\). Since
\[
 |(x+iy)^2+t^2|\ge x^2+a_0^2,
\]
the nonzero Fourier coefficients are bounded by
\(C(s)a_0^{1-2s}e^{-2\pi n(t-a_0)}\). Summing both signs and
all \(b,n\ge1\), put
\[
 z=e^{-2\pi v_-},\quad y_0=e^{-2\pi(v_--a_0)},\qquad
 E_A=\frac{2C(s)a_0^{1-2s}y_0}{(1-z)(1-y_0)}.
\]
Then \(|\mathcal A_q(s)-\mathcal F_q(s)|\le E_A\) uniformly in
the band.

For the \(s\)-derivative, the complex argument of
\((x+iy)^2+t^2\) has absolute value at most \(\pi/2\), and its
logarithm is bounded in modulus by
\[
 \log(1+x^2/a_0^2)+2\log(2t)+2|\log a_0|+\pi/2.
\]
The integral of the first term against
\((x^2+a_0^2)^{-s}\) is
\(C(s)a_0^{1-2s}g(s)\). Using \(\log b\le b-1\) in the remaining
geometric sum therefore gives
\begin{equation}\label{EN}
 |\mathcal A_q'(s)-\mathcal F_q'(s)|\le E_N,
 \quad E_N=E_A\left\{2\log(2v_+)+2|\log a_0|+g(s)+\frac\pi2
                                      +\frac{2z}{1-z}\right\}.
\end{equation}
The contour shift is justified by decay at infinity for \(s>1\);
the same decay with an extra logarithm justifies differentiation.

Put \(d=s-1\), \(z_2=\zeta(2s)\),
\(c_s=-2\zeta'(2s)/\zeta(2s)\), and
\[
 \kappa_d=g(s)-2\frac{\zeta'(2s-1)}{\zeta(2s-1)}-c_s.
\]
The main upper term is
\begin{equation}\label{Hmain}
 H_q(d)=\frac{\mathcal F_q(s)\{\log(q/4)+\kappa_d\}}
                   {2\{z_2+\mathcal F_q(s)\}}.
\end{equation}
The certificate checks that \(H_q(d)>0\) and decreases with \(q\)
for each exponent used. Since \(z_2>1\), the logarithmic-derivative
error is at most
\begin{equation}\label{Herror}
 E_H(d)=\frac{E_N+(c_s+2H_{Q_0}(d))E_A}{2(1-E_A)}.
\end{equation}
Consequently a uniform upper bound for the weighted noninert sum is
\begin{equation}\label{weightedupper}
 U(q)=\sum_i c_iF(d_i)H_q(d_i)
       +\sum_i|c_i|F(d_i)E_H(d_i)
       +\frac{\log Q_0}{2(Q_0-1)}.
\end{equation}
The error is bounded after all coefficient magnitudes are included.

\section{Exact certification and completion of the proof}\label{certificate}
The certificate uses the five bands
\[
 [570,580],\quad[580,600],\quad[600,700],\quad[700,1050],\quad[1050,2000].
\]
The first uses a polynomial of degree 16 with \(d=3/10,e=1/25\);
the others use degree 12 with \(d=2/5,e=1/20\). The integer
polynomial coefficients, auxiliary-prime rows, and eight-stage
refinement schedule are fixed in \texttt{parameters.json}.

For completeness, the dependence on \(q\) is certified over each
entire band. Write \(Y=\log(q/Q_0)\),
\(a_i=1/2+d_i\), and
\(r_i=\mathcal F_{Q_0}(s_i)/\zeta(2s_i)\). Expand
\[
 H_q(d_i)=\frac12(\log(Q_0/4)+Y+\kappa_{d_i})
       \sum_{\ell\ge1}(-1)^{\ell+1}r_i^\ell e^{-\ell a_iY}.
\]
The certificate bounds the differentiated geometric tail explicitly.
For each retained \(\ell\), the signed finite sum is evaluated as
a polynomial in \(1-e^{-\ell eY}\), with its coefficients computed
by exact interval arithmetic. This avoids naive interval cancellation.
Cells of width \(1/1000\) cover the entire \(Y\)-interval. On every
cell the derivative of the main sum in \eqref{weightedupper} is
strictly less than \(-\gamma\). The two error bounds in that equation
are uniform constants, so they require no differentiation.

\IfFileExists{numerical_table.tex}{\begin{center}
\small
\begin{tabular}{@{}rrrrr@{}}
\toprule
Band & Lower bound & Upper bound & Gap $>$ & Norm error $<$\\
\midrule
$[570,580]$ & 10.975413 & 10.965024 & 0.010389 & $7.652\times10^{-5}$\\
$[580,600]$ & 9.566392 & 9.544454 & 0.021938 & $1.969\times10^{-10}$\\
$[600,700]$ & 9.755174 & 9.682808 & 0.072367 & $4.675\times10^{-11}$\\
$[700,1050]$ & 9.890348 & 9.507108 & 0.383240 & $6.739\times10^{-14}$\\
$[1050,2000]$ & 12.460477 & 11.013702 & 1.446776 & $3.672\times10^{-22}$\\
\bottomrule
\end{tabular}
\end{center}
All displayed bounds are rounded outward. The upper bound includes
both the norm-series error and the ramified-prime allowance.
The exact certificate endpoints are recorded in the JSON output.
}{}

At \(q=Q_0\), the exact lower endpoint of
\eqref{weightedlower} is strictly greater than the exact upper
endpoint of \eqref{weightedupper}. Since the latter decreases faster
than \(-\gamma\log(q/Q_0)\), \eqref{Luniform} proves a contradiction
throughout the band. The five bands exclude \([570,2000]\); the
audited v5 implication excludes \(q\ge2000\), proving
Theorem~\ref{main}.

The numerical proof uses integer and rational intervals throughout.
Logarithms use a convergent positive series after range reduction;
exponentials use a Taylor series with an explicit tail; square roots
use integer square roots. Zeta values and derivatives use
Euler--Maclaurin with the periodic-Bernoulli remainder, while gamma
and digamma use positive-real Stirling expansions with remainder
bounds. These are evaluations in \(\Re s>1\), not zero-location
inputs. Acceptance never depends on a floating-point optimizer.

Run, without Python's \texttt{-O} option,
\begin{verbatim}
python3 verify_v6.py
python3 verify_structures.py
python3 verify_finite.py
\end{verbatim}
The first command certifies the inequalities and writes exact
endpoints to JSON. The second checks finite instances of the new
algebraic identities and divisibility statements; these tests do not
replace their proofs. The third is a separate finite check below the
analytic cutoff and is not used to establish that cutoff.

\subsection{Numerical remainder bounds}
Here are the remainder formulas used in the certificate; the integer
parameters in this paragraph are numerical evaluation parameters,
independent of the determinant dimensions.
For \(s>1\), Euler--Maclaurin with \(N_0=64\), \(m=18\) gives
\begin{align*}
 \zeta(s)={}&\sum_{n<N_0}n^{-s}+\frac{N_0^{1-s}}{s-1}
      +\frac12N_0^{-s}
      +\sum_{k=1}^{m}\frac{B_{2k}(s)_{2k-1}}{(2k)!}
                      N_0^{1-s-2k}+\mathcal R,\\
 |\mathcal R|\le{}&E=
 \frac{|B_{2m}|(s)_{2m}}{(2m)!(s+2m-1)}N_0^{1-s-2m},\\
 |\mathcal R'|\le{}&E\left\{\log N_0+
       \sum_{j=0}^{2m-1}\frac1{s+j}+\frac1{s+2m-1}\right\}.
\end{align*}
Here \((s)_j\) denotes the rising factorial. These bounds follow
from the periodic-Bernoulli integral and its derivative.
For log gamma and digamma, shift to \(z=x+64\), retain the terms
\(1\le k<20\) in the positive-real Stirling expansions, and use
the first omitted term as an absolute remainder bound
\cite{dlmf}:
\[
 |R_{\log\Gamma}|\le\frac{|B_{40}|}{40\cdot39\,z^{39}},
 \qquad |R_\psi|\le\frac{|B_{40}|}{40z^{40}}.
\]
The recurrence formulas then recover the values at \(x\).

After range reduction to \(1\le y\le2\), logarithms use 80 terms
of \(2\sum_{k\ge0}t^{2k+1}/(2k+1)\), where
\(t=(y-1)/(y+1)\), and tail
\(2t^{161}/(161(1-t^2))\). Exponentials are reduced to
\(0\le x\le1\); the series through degree 80 has positive tail
at most \(82/(81\cdot81!)\). Repeated squaring and reciprocation
undo the reduction. Machin's formula with 60 arctangent terms
encloses \(\pi\). Elementary interval operations round outward
to 65 decimal places; square roots use integer bounds.

For the differentiated geometric series above, the tail after
\(L\) terms is at most
\[
 \frac12\sum_i |c_i|F(d_i)r_i^{L+1}
 \left\{\frac1{1-r_i}+a_i E_*
       \left(\frac{L+1}{1-r_i}+\frac{r_i}{(1-r_i)^2}\right)\right\},
 \qquad E_*=\max_i\{\log(Q_1/4)+\kappa_{d_i}\}.
\]
The certificate increases \(L\) until this bound is below
\(10^{-10}\), then includes it in every derivative cell test.

\section{Directions tested and the stopping point}
The research continued beyond the first successful lighter-method
bound. The retained improvements include disk interpolation, the
prime itself as monomial width, a mixed ideal sieve, higher Frobenius
factors, the smoothed valuation sum, direct finite factorial refinement,
the joint radius/order inequality, and optimized positive kernels.

Other routes were examined but are not inputs to this proof.
Weighted monomial regions suggested by Kerr lose the exact packing
of the degrees \(r+p_0s\) used here. Numerical tests with disk
orthogonal polynomials gave only a small further improvement over
the joint monomial estimate. Finer transport bounds and more elaborate
kernels offered diminishing returns, with increasingly difficult
signed-sum error control. The sieve work of Gaudet \cite{gaudet}
was also considered; its additional class-group analytic inputs and
unspecified numerical exponent did not supply a useful replacement
for the present explicit argument. The quartic-field/Masser and
quadratic zero-free approaches were not imported.

The cutoff 570 is a practical stopping point for this research pass,
not an optimality theorem about these methods. The package records
the final proof and certificate as one v6 result, and preserves the
distinction between exact arithmetic, mathematical audit, and formal
verification.

\begin{thebibliography}{99}
\bibitem{tao} T.~Tao, \emph{An experimental AI-math challenge: using
tenth-discriminant bounds to explore arguments ruling out Siegel zeroes},
20 October 2026.
\href{https://terrytao.wordpress.com/2026/10/20/an-experimental-ai-math-challenge-using-tenth-discriminant-bounds-to-explore-arguments-ruling-out-siegel-zeroes/}{Challenge and discussion}.
\bibitem{taocomment} T.~Tao, comment 698471 on \cite{tao}.
\href{https://terrytao.wordpress.com/2026/10/20/an-experimental-ai-math-challenge-using-tenth-discriminant-bounds-to-explore-arguments-ruling-out-siegel-zeroes/\#comment-698471}{Request for bounds using lighter machinery}.
\bibitem{green} B.~Green, \emph{On Landau--Siegel zeros}, v2, 2026.
\href{https://hexagonmath.org/2610.00311v2}{hexagon:2610.00311v2}.
Its provenance statement credits the mathematical arguments to ChatGPT
and the exposition to Green.
\bibitem{vandijk} R.~van Dijk, \emph{A Frobenius determinant proof of
the non-existence of Landau--Siegel zeros}, v1, 2026.
\href{https://hexagonmath.org/2610.00409v1}{hexagon:2610.00409v1}.
The finite-field argument required here is proved in Section~\ref{rank}.
\bibitem{greencomment} B.~Green, comment 698463 on \cite{tao}.
\href{https://terrytao.wordpress.com/2026/10/20/an-experimental-ai-math-challenge-using-tenth-discriminant-bounds-to-explore-arguments-ruling-out-siegel-zeroes/\#comment-698463}{Auxiliary-prime rank argument and Pescador attribution}.
\bibitem{pratt} K.~Pratt, comment 698476 on \cite{tao}.
\href{https://terrytao.wordpress.com/2026/10/20/an-experimental-ai-math-challenge-using-tenth-discriminant-bounds-to-explore-arguments-ruling-out-siegel-zeroes/\#comment-698476}{Norm counts, stronger divisibility, and disk interpolation}.
\bibitem{kerrcomment} B.~Kerr, comment 698475 on \cite{tao}.
\href{https://terrytao.wordpress.com/2026/10/20/an-experimental-ai-math-challenge-using-tenth-discriminant-bounds-to-explore-arguments-ruling-out-siegel-zeroes/\#comment-698475}{Monomial regions and higher Frobenius relations}.
See also \emph{A simpler proof that there are no Landau--Siegel zeros},
\href{https://arxiv.org/abs/2610.21099}{arXiv:2610.21099}.
\bibitem{ross} michaelmross, comment 698466 on \cite{tao}.
\href{https://terrytao.wordpress.com/2026/10/20/an-experimental-ai-math-challenge-using-tenth-discriminant-bounds-to-explore-arguments-ruling-out-siegel-zeroes/\#comment-698466}{Poisson expansion of the principal quadratic form}.
\bibitem{rs} J.~Barkley Rosser and L.~Schoenfeld,
\emph{Approximate formulas for some functions of prime numbers},
Illinois J. Math. \textbf{6} (1962), 64--94.
\href{https://doi.org/10.1215/ijm/1255631807}{doi:10.1215/ijm/1255631807}.
Theorems 4, 6, 9 and 14, and Table I on printed page 90.
\bibitem{v5} \emph{A Mellin-averaged determinant bound for the tenth
discriminant}, version 0.5.0.
\href{https://github.com/twinheterodox/tenth-discriminant/releases/tag/v0.5.0}{Published v5 release},
including its minor Section 6.2 wording correction. The audited
implication \(h(K)=1\Rightarrow q<2000\) is the earlier-version input.
The corrected source is pinned at
\href{https://github.com/twinheterodox/tenth-discriminant/tree/6f781492afd06de32997c2c89d4e228a78bab7fc/versions/v5}{commit 6f781492}.
\bibitem{gaudet} L.~M.~Gaudet, \emph{Sifting for small split primes of
an imaginary quadratic field in a given ideal class},
\href{https://arxiv.org/abs/2408.01610}{arXiv:2408.01610}.
\bibitem{dlmf} NIST Digital Library of Mathematical Functions,
\href{https://dlmf.nist.gov/2.10}{Section 2.10 (Euler--Maclaurin)} and
\href{https://dlmf.nist.gov/5.11}{Section 5.11 (Stirling expansions
and positive-real remainder bounds)}.
\end{thebibliography}
\end{document}
